% Standalone problem document, generated from the adjacent README.md.
% Compile with XeLaTeX. Mathematical content is maintained in Markdown.
\documentclass[11pt,a4paper]{article}
\usepackage[fontsize=10.5pt]{scrextend}
\usepackage[margin=22mm,headheight=15pt,headsep=9mm,footskip=11mm]{geometry}
\usepackage{fontspec}
\setmainfont{texgyrepagella}[Extension=.otf,UprightFont=*-regular,BoldFont=*-bold,ItalicFont=*-italic,BoldItalicFont=*-bolditalic]
\setsansfont{texgyreheros}[Extension=.otf,UprightFont=*-regular,BoldFont=*-bold,ItalicFont=*-italic,BoldItalicFont=*-bolditalic]
\setmonofont{texgyrecursor}[Extension=.otf,UprightFont=*-regular,BoldFont=*-bold,ItalicFont=*-italic,BoldItalicFont=*-bolditalic,Scale=MatchLowercase]
\usepackage{amsmath,amssymb}
\usepackage{unicode-math}
\setmathfont{texgyrepagella-math.otf}
\usepackage{xcolor,microtype,enumitem,needspace,fancyhdr}
\usepackage{bookmark}
\definecolor{ink}{HTML}{17344B}
\definecolor{accent}{HTML}{197D80}
\definecolor{muted}{HTML}{596773}
\hypersetup{colorlinks=true,linkcolor=accent,urlcolor=accent,pdftitle={TR-21 - A
Seginer theorem for arbitrary independent identically distributed tensor
entries},pdfauthor={Open Problems in Numerical Linear Algebra}}
\urlstyle{same}
\setlength{\parindent}{0pt}
\setlength{\parskip}{5pt plus 1pt minus 1pt}
\linespread{1.04}
\setlength{\emergencystretch}{3em}
\widowpenalty=10000
\clubpenalty=10000
\displaywidowpenalty=10000
\allowdisplaybreaks[1]
\setlist{nosep,leftmargin=1.5em}
\providecommand{\tightlist}{\setlength{\itemsep}{0pt}\setlength{\parskip}{0pt}}
\providecommand{\pandocbounded}[1]{#1}
\pagestyle{fancy}
\fancyhf{}
\fancyhead[L]{\sffamily\footnotesize\color{muted}OPEN PROBLEMS IN NUMERICAL LINEAR ALGEBRA}
\fancyhead[R]{\sffamily\bfseries\footnotesize\color{accent}TR-21}
\fancyfoot[L]{\parbox[b]{0.9\textwidth}{\sffamily\fontsize{8}{10}\selectfont\color{muted}Editorial ratings; a bounded search cannot certify that a problem remains unsolved.\\Literature check: 2026-09-10}}
\fancyfoot[R]{\sffamily\footnotesize\color{muted}\thepage}
\renewcommand{\headrulewidth}{0.3pt}
\renewcommand{\headrule}{\hbox to\headwidth{\color{accent}\leaders\hrule height \headrulewidth\hfill}}
\makeatletter
\renewcommand{\section}{\Needspace{5\baselineskip}\@startsection{section}{1}{0pt}{12pt plus 2pt minus 2pt}{4pt}{\sffamily\bfseries\large\color{ink}}}
\renewcommand{\subsection}{\Needspace{4\baselineskip}\@startsection{subsection}{2}{0pt}{9pt plus 1pt minus 1pt}{3pt}{\sffamily\bfseries\normalsize\color{ink}}}
\makeatother
\setcounter{secnumdepth}{0}
\begin{document}
{\sffamily\small\color{muted}Tensor computations\par}
\vspace{3pt}
{\raggedright\hyphenpenalty=10000\exhyphenpenalty=10000\sffamily\bfseries\fontsize{21}{24}\selectfont\color{ink}A
Seginer theorem for arbitrary independent identically distributed tensor
entries\par}
\vspace{7pt}
{\sffamily\small\textbf{Difficulty:} challenging\quad\textbf{Importance:} interesting
to the community\par}
{\sffamily\small\bfseries\color{accent}Status: Partially resolved\par}
\vspace{4pt}
{\color{accent}\hrule height 0.8pt}
\vspace{6pt}
\textbf{Rating rationale:} Obtaining a uniform bound without moment
control needs substantial new random-tensor analysis. Sharp
injective-norm estimates matter to the community working on tensor
approximation and random multilinear algorithms.

\subsection{Problem statement}\label{problem-statement}

Fix an integer \(r\geq 3\). For integers \(n_1,\ldots,n_r\geq2\), let
\(T\in\mathbb R^{n_1\times\cdots\times n_r}\) have independent
identically distributed integrable real entries of mean zero. Its
injective norm is \[
\|T\|_{\mathrm{inj}}=
\sup_{\|x_j\|_2=1,\ 1\leq j\leq r}
|\langle T,x_1\otimes\cdots\otimes x_r\rangle|.
\] Define the largest expected Euclidean fiber norm by \[
F(T)=\max_{1\leq k\leq r}\mathbb E
\max_{\substack{i_j\in\{1,\ldots,n_j\}\\j\ne k}}
\left(\sum_{a=1}^{n_k}
T_{i_1,\ldots,i_{k-1},a,i_{k+1},\ldots,i_r}^{\,2}\right)^{1/2}.
\] Do constants \(0<c_r\leq C_r<\infty\), depending only on the order
\(r\), exist such that \[
c_rF(T)\leq\mathbb E\|T\|_{\mathrm{inj}}\leq C_rF(T)
\] for every such format and entry distribution?

This is Lucca--Pesenti's Conjecture 1.8, with its comparison notation
written as explicit quantifiers. The entry distribution may depend on
the dimensions. No variance normalization, finite higher moment,
density, or dimension-independent moment-equivalence hypothesis is
imposed. Integrability makes both expectations finite in every finite
format. The matrix case \(r=2\) is Seginer's theorem and is included
only as background.

\subsection{Why it matters}\label{why-it-matters}

The injective norm measures the largest interaction with a rank-one
tensor. A fiber-based characterization would give sharp noise scales for
random tensor approximation, including sparse distributions for which
flattening and logarithmic losses can obscure the relevant scale.

\newpage
\subsection{References and status
check}\label{references-and-status-check}

\begin{enumerate}
\def\labelenumi{\arabic{enumi}.}
\tightlist
\item
  K. Lucca and L. Pesenti, \emph{Norm Bounds for Sparse Random Tensors
  and Spectral Gap of Random Hypergraphs},
  \href{https://arxiv.org/html/2607.07308v1}{arXiv:2607.07308v1}, §1.3,
  Conjecture 1.8; Assumption 1.6, Theorem 1.7, and Appendix B describe
  the proved moment-restricted case.
\item
  Y. Seginer, \emph{The expected norm of random matrices},
  Combinatorics, Probability and Computing \textbf{9} (2000), 149--166;
  the matrix predecessor, identified in reference 26 of Lucca--Pesenti.
\end{enumerate}

\subsubsection{Status check ---
2026-09-10}\label{status-check-2026-09-10}

Checked Lucca--Pesenti v1, Conjecture 1.8, Theorem 1.7 and Proposition
B.1, and targeted Seginer/tensor/2026 resolution searches. Their
moment-restricted theorem proves the comparison for subclasses with
fixed numerical moment-equivalence parameters; its constants otherwise
depend on those parameters as well as the order. These are substantive
proved subclasses, so the status is Partially resolved. The requested
uniformity over arbitrary integrable distributions remains open,
explicitly including centered Bernoulli entries with parameter \(1/n\).
No full resolution was located. The matrix theorem and estimates with
logarithmic losses do not settle the displayed target.
\end{document}
